\subsection{\texorpdfstring{The $U(1)$ Problem}{The U(1) Problem}}
\label{sec:19.10}

The success of quantum chromodynamics in explaining the pattern of strong-interaction symmetries seemed at first to be marred by one failure. As we saw in Section 19.5, the broken $SU(2)\times SU(2)$ symmetry of the strong interactions is a natural consequence of the smallness of the $u$ and $d$ quark masses. But the Lagrangian of quantum chromodynamics with small $u$ and $d$ quark masses also has another chiral symmetry~\hyperref[ch19-ref-51]{[51]}, the $U(1)_A$ symmetry under transformations
\begin{equation}
u\longrightarrow\exp(\ii\gamma_5\theta)u,
\qquad
\dd\longrightarrow\exp(\ii\gamma_5\theta)d.
\tag{19.10.1}\label{eq:19.10.1}
\end{equation}
Such a symmetry if unbroken would, like $SU(2)\times SU(2)$, impose a parity doubling on the hadron spectrum, but no such parity doubling is observed. On the other hand, a broken $U(1)_A$ symmetry would imply the existence of an isoscalar $0^-$ Goldstone boson with a mass comparable to that of the pion, and this Goldstone boson is also not observed. It is true that the $\eta$ meson is an isoscalar $0^-$ boson, but it is considerably heavier than the pion, and as we saw in Section 19.7, it is well understood as one of the Goldstone bosons of $SU(3)\times SU(3)$. With the $s$ as well as the $u$ and $d$ quarks regarded as relatively light, the Lagrangian of quantum chromodynamics would have a $U(1)$ chiral symmetry in addition to $SU(3)\times SU(3)$, under the transformations
\begin{equation}
u\longrightarrow\exp(\ii\gamma_5\theta)u,
\qquad
\dd\longrightarrow\exp(\ii\gamma_5\theta)d,
\qquad
s\longrightarrow\exp(\ii\gamma_5\theta)s.
\tag{19.10.2}\label{eq:19.10.2}
\end{equation}
The spontaneous breakdown of such a symmetry would require the existence of \emph{two} isoscalar $0^-$ mesons: one the $\eta$, and the other with a mass comparable to that of the pion.

This prediction can be made more explicit~\hyperref[ch19-ref-52]{[52]}. If we include the transformation \eqref{eq:19.10.2} among the spontaneously broken symmetries of massless quantum chromodynamics, then in the presence of quark masses we encounter a term \eqref{eq:19.7.12}:
\begin{equation}
\mathcal L_{\mathrm{mass}}
=-\overline q M_q q
=-\overline{\widetilde q}\,
e^{-\ii\sqrt{2}\gamma_5B/F_\pi}
M_q
e^{-\ii\sqrt{2}\gamma_5B/F_\pi}
\widetilde q,
\tag{19.10.3}\label{eq:19.10.3}
\end{equation}
where again
\begin{equation}
M_q=
\begin{bmatrix}
m_u&0&0\\
0&m_d&0\\
0&0&m_s
\end{bmatrix},
\tag{19.10.4}\label{eq:19.10.4}
\end{equation}
but now $B$ includes a Goldstone boson field $\zeta$ for the broken symmetry \eqref{eq:19.10.2}:
\begin{align}
B={}&
\begin{bmatrix}
\dfrac{1}{\sqrt2}\pi^0+\dfrac{1}{\sqrt6}\eta^0
&\pi^+&K^+\\[4pt]
\pi^-&
-\dfrac{1}{\sqrt2}\pi^0+\dfrac{1}{\sqrt6}\eta^0
&K^0\\[4pt]
\overline K^-&\overline K^0&-\sqrt{\dfrac23}\eta^0
\end{bmatrix}
\notag\\
&+\frac{F_\pi}{\sqrt3F_\zeta}
\begin{bmatrix}
\zeta&0&0\\
0&\zeta&0\\
0&0&\zeta
\end{bmatrix},
\tag{19.10.5}\label{eq:19.10.5}
\end{align}
with $F_\zeta$ the unknown coupling of the $U(1)_A$ Goldstone boson to the corresponding current (and the factor $\sqrt3$ inserted for future convenience.) We can again use the unbroken $SU(3)$ symmetry to write the vacuum expectation value of the quark bilinear in the form \eqref{eq:19.7.14}, and expanding in powers of the boson fields, we find a Goldstone boson mass term in the Lagrangian of the form
\begin{align}
-\frac{v}{F_\pi^2}
\operatorname{Tr}\bigl\{B,\{B,M_q\}\bigr\}
={}&-\frac{v}{F_\pi^2}\Bigg[
4m_u\left(
\frac{1}{\sqrt2}\pi^0+\frac{1}{\sqrt6}\eta^0
+\frac{F_\pi}{\sqrt3F_\zeta}\zeta
\right)^2
\notag\\
&\quad+4(m_u+m_d)\pi^+\pi^-
+4(m_u+m_s)K^+\overline K^-
\notag\\
&\quad+4m_d\left(
-\frac{1}{\sqrt2}\pi^0+\frac{1}{\sqrt6}\eta^0
+\frac{F_\pi}{\sqrt3F_\zeta}\zeta
\right)^2
\notag\\
&\quad+4(m_d+m_s)K^0\overline K^0
\notag\\
&\quad+4m_s\left(
-\sqrt{\frac23}\eta^0
+\frac{F_\pi}{\sqrt3F_\zeta}\zeta
\right)^2
\Bigg].
\tag{19.10.6}\label{eq:19.10.6}
\end{align}
The charged and strange meson masses are the same as before, but now the neutral non-strange mesons have a mass matrix
\begin{equation}
M_0^2=8v
\begin{bmatrix}
\dfrac{m_u+m_d}{2F_\pi^2}
&
\dfrac{m_u-m_d}{2\sqrt3F_\pi^2}
&
\dfrac{m_u-m_d}{\sqrt6F_\pi F_\zeta}
\\[8pt]
\dfrac{m_u-m_d}{2\sqrt3F_\pi^2}
&
\dfrac{m_u+m_d+4m_s}{6F_\pi^2}
&
\dfrac{m_u+m_d-2m_s}{3\sqrt2F_\pi F_\zeta}
\\[8pt]
\dfrac{m_u-m_d}{\sqrt6F_\pi F_\zeta}
&
\dfrac{m_u+m_d-2m_s}{3\sqrt2F_\pi F_\zeta}
&
\dfrac{m_u+m_d+m_s}{3F_\zeta^2}
\end{bmatrix}
\tag{19.10.7}\label{eq:19.10.7}
\end{equation}
with rows and columns listed in the order $\pi^0$, $\eta^0$, and $\zeta$.

In the limit where $m_u$ and $m_d$ vanish, this has two eigenvectors with eigenvalue zero:
\begin{equation}
u_a=
\begin{pmatrix}1\\0\\0\end{pmatrix},
\qquad
u_b=
\frac{1}{\sqrt{F_\pi^2+2F_\zeta^2}}
\begin{pmatrix}0\\F_\pi\\\sqrt2F_\zeta\end{pmatrix}.
\tag{19.10.8}\label{eq:19.10.8}
\end{equation}
In this orthonormal basis, the mass-squared matrix to first order in $m_u$ and $m_d$ is
\begin{equation}
m_{aa}^2\simeq u_a^{\mathsf T}M_0^2u_a
=\frac{4v(m_u+m_d)}{F_\pi^2},
\tag{19.10.9}\label{eq:19.10.9}
\end{equation}
\begin{equation}
m_{bb}^2\simeq u_b^{\mathsf T}M_0^2u_b
=\frac{12v(m_u+m_d)}{F_\pi^2+2F_\zeta^2},
\tag{19.10.10}\label{eq:19.10.10}
\end{equation}
\begin{equation}
m_{ab}^2=m_{ba}^2
\simeq
\frac{\sqrt3v(m_u-m_d)}
{2F_\pi\sqrt{F_\pi^2+2F_\zeta^2}}.
\tag{19.10.11}\label{eq:19.10.11}
\end{equation}

The effect of the off-diagonal term \eqref{eq:19.10.11} is to decrease the product of eigenvalues by a negligible fractional amount $(m_u-m_d)^2/64(m_u+m_d)^2$, while of course leaving the sum of the eigenvalues unchanged, so the eigenvalues are given to a good approximation by the diagonal elements \eqref{eq:19.10.9} and \eqref{eq:19.10.10}. Comparing Eq.~\eqref{eq:19.10.9} with Eq.~\eqref{eq:19.7.16}, we see that the particle corresponding to the eigenvector $u_a$ is the $\pi^0$. The other particle, corresponding to $u_b$, has a mass
\begin{equation}
m_b\simeq\sqrt{m_{bb}^2}
\simeq
\frac{\sqrt3m_\pi F_\pi}
{\sqrt{F_\pi^2+2F_\zeta^2}}
\leq\sqrt3m_\pi.
\tag{19.10.12}\label{eq:19.10.12}
\end{equation}
Thus a broken $U(1)_A$ symmetry would require a neutral pseudoscalar Goldstone boson with mass less than $\sqrt3m_\pi$, in addition to the pion itself. It hardly needs to be said that no such strongly-interacting particle exists.\footnote{Note that taking $F_\zeta\gg F_\pi$ would give the extra neutral scalar a very light mass, but its interactions would be so weak that it might have escaped detection. It seems unlikely that a very large ratio of $F_\pi$ to $F_\zeta$ could arise in quantum chromodynamics, but something like this happens in theories with extra field variables that have been proposed to avoid parity violation by instantons; see Section 23.6.} We shall see in Section 23.5 that this problem was eventually solved by the discovery of non-perturbative effects that violate the extra $U(1)_A$ symmetry.
